%%%%%%%%%%%%
%
% $Autor: Wings $
% $Datum: 2019-03-05 08:03:15Z $
% $Pfad: Sketch.tex $
% $Version: 4250 $
% !TeX spellcheck = en_GB/de_DE
% !TeX encoding = utf8
% !TeX root = manual 
% !TeX TXS-program:bibliography = txs:///biber
%
%%%%%%%%%%%%

\chapter{Sketch}

\section{System Sketch}

The following sketch illustrates the overall hardware architecture and data flow of the \textbf{License Plate Detection System}.  
It represents how each physical component connects to the Jetson Nano during operation — from image capture to output display and database logging.  
This configuration shows the system as it would appear during actual deployment in a parking or roadside monitoring setup.

% --------------------------------------------------
% TikZ Hardware Layout Diagram
% --------------------------------------------------
\begin{figure}[h!]
	\centering
	\resizebox{1.1\textwidth}{!}{% scales safely within margins
		\begin{tikzpicture}[node distance=2.6cm, >=stealth, thick]
			\tikzstyle{block}=[rectangle, draw=blue!60!black, fill=blue!10,
			rounded corners, text centered, text width=3.6cm, minimum height=1.2cm]
			\tikzstyle{db}=[cylinder, shape border rotate=90, draw=orange!80!black,
			fill=orange!10, minimum height=1.3cm, minimum width=1cm, aspect=0.4, align=center]
			
			% Nodes
			\node[block] (camera) {\faCamera \\ \textbf{USB Camera}\\(Microsoft LifeCam Show)};
			\node[block, right=2.9cm of camera] (nano) {\faMicrochip \\ \textbf{NVIDIA Jetson Nano}\\(reComputer J1010)};
			\node[block, right=2.9cm of nano] (display) {\faDesktop \\ \textbf{HDMI Display}\\(Real-time Output)};
			\node[db, below=1.7cm of nano] (db) {\textbf{SQLite Database}\\(plates.db)};
			\node[block, below=1.7cm of camera] (power) {\faBolt \\ \textbf{Power Supply}\\(5 V / 4 A USB-C)};
			
			% Connections
			\draw[->, thick] (camera) -- node[above, align=center, font=\scriptsize]{Video Stream\\(USB 2.0)} (nano);
			\draw[->, thick] (nano) -- node[above, align=center, font=\scriptsize]{Output\\(HDMI)} (display);
			\draw[->, thick] (nano) -- node[right, align=center, font=\scriptsize]{(SQL)} (db);
			\draw[->, thick] (power) -- node[left, align=center, font=\scriptsize]{Power Input} (nano);
			
		\end{tikzpicture}
	}% end resizebox
	\caption{Functional hardware sketch of the Jetson Nano License Plate Detection System}
	\label{fig:hardware_sketch}
\end{figure}

\vspace{0.3cm}

\noindent
The hardware layout above shows the system’s logical and physical connections.  
The USB camera continuously streams live video to the Jetson Nano for inference.  
The HDMI port drives a real-time display showing detection output, while the database logging occurs in the background.  
The unit is powered by a 5V/4A supply through a USB-C input.  
This modular structure allows easy assembly, transport, and maintenance in embedded AI environments.

% --------------------------------------------------
% New Image: Jetson Nano Aluminum Case with Ports
% --------------------------------------------------
\begin{figure}[h!]
	\centering
	\includegraphics[width=0.72\textwidth]{General/jetsonNanoPorts.png}
	\caption{Aluminum-encased Jetson Nano (Seeed Studio reComputer J1010) showing external ports for HDMI, USB, Ethernet, and power.}
	\label{fig:jetson_case_ports}
\end{figure}

\noindent
The above image shows the actual Seeed Studio reComputer J1010 aluminum enclosure used in this project.  
The front and side ports include:
\begin{itemize}
	\item Two \textbf{USB 3.0} ports for peripherals such as cameras or flash drives.
	\item One \textbf{HDMI} output port for display connection.
	\item One \textbf{Gigabit Ethernet} port for optional network connectivity.
	\item One \textbf{USB-C} input for 5V/4A power supply.
\end{itemize}

\noindent
The compact design and passive cooling of the aluminum casing allow the system to operate continuously under load while maintaining thermal stability.  
This makes the setup suitable for continuous field operation, even in confined or semi-outdoor conditions.
\vspace{0.3cm}
\noindent
\textbf{Practical Deployment Notes:}  
The reComputer J1010 unit provides a highly integrated, plug-and-play platform for real-world AI applications.  
Because of its small form factor (dimensions: \textbf{130mm × 120mm × 45mm}), it can be easily mounted beside surveillance cameras, on dashboards, or inside parking enclosures.  
The fanless aluminum casing not only reduces noise but also enhances durability and heat dissipation, which is essential for continuous inference tasks such as YOLOv8-based object detection.

\begin{table}[h!]
	\centering
	\caption{Jetson Nano Hardware and Connectivity Summary}
	\label{tab:jetson_summary}
	\renewcommand{\arraystretch}{1.3} % increases row spacing
	\begin{tabular}{|l|l|}
		\hline
		\textbf{Feature} & \textbf{Description} \\
		\hline
		Processor & Quad-core ARM Cortex-A57 CPU \\
		\hline
		GPU & 128-core Maxwell architecture \\
		\hline
		Memory & 4 GB LPDDR4 \\
		\hline
		Operating System & JetPack 5.1.2 (Ubuntu 20.04) \\
		\hline
		Camera Interface & USB 2.0 (Microsoft LifeCam Show) \\
		\hline
		Display Interface & HDMI 2.0, 4K capable output \\
		\hline
		Networking & Gigabit Ethernet LAN port \\
		\hline
		Storage & 64 GB microSD card \\
		\hline
		Power Input & 5V / 4A via USB-C \\
		\hline
		Cooling System & Passive aluminum heat sink (fanless) \\
		\hline
	\end{tabular}
\end{table}

\vspace{0.3cm}
\noindent
The hardware configuration summarized in Table~\ref{tab:jetson_summary} ensures that the Jetson Nano can perform \textbf{real-time object detection and OCR processing} while maintaining a low thermal footprint.  
Its low power consumption and compact design make it an ideal choice for embedded edge-AI applications such as traffic surveillance, parking automation, and smart tolling systems.

\noindent
In this project, all computations — including YOLOv8 inference, EasyOCR text recognition, and SQLite database logging — are executed locally on the Jetson Nano without any dependency on cloud services.  
This ensures complete \textbf{data privacy, low latency, and offline operability}, which are key requirements for on-premise vehicle monitoring deployments.
\vspace{0.4cm}
\noindent
\textbf{System Integration and Workflow:}  
The hardware design of the License Plate Detection System is optimized for seamless integration with the accompanying software stack.  
The Jetson Nano acts as both the computation and control hub, handling camera input, model inference, and result logging in a single environment.  
This all-in-one configuration minimizes latency between hardware operations and ensures stable frame processing without the need for external servers.

\noindent
From a software perspective, the Jetson Nano runs a lightweight Ubuntu environment based on the \textbf{JetPack 5.1.2 SDK}, which provides direct support for CUDA, cuDNN, TensorRT, and OpenCV libraries.  
These NVIDIA components enable accelerated deep-learning inference directly on the device GPU, providing high throughput for YOLOv8 detection even under edge constraints.

\noindent
The following points summarize the advantages of this embedded integration:
\begin{itemize}
	\item \textbf{Low Latency:} On-device computation eliminates round-trip delay associated with cloud processing.  
	\item \textbf{Compact Footprint:} The hardware, camera, and power supply together occupy less than 0.5 L of volume, making it deployable in small enclosures.  
	\item \textbf{Energy Efficiency:} Despite continuous GPU usage, the system remains under 10 W average power draw.  
	\item \textbf{Reliability:} Solid aluminum casing and passive cooling ensure uninterrupted operation over extended periods.  
	\item \textbf{Scalability:} Multiple Jetson units can operate in parallel for large-scale installations such as parking lots or toll plazas.  
\end{itemize}

\noindent
Overall, this hardware architecture provides an ideal foundation for embedded computer-vision applications.  
Its modular design allows direct adaptation for related use cases such as \textit{traffic violation detection}, \textit{vehicle counting}, or \textit{entry-exit automation}.  
